\centerline{\bf \Huge Разрезания на доминошки}

\begin{flushright}
        \scriptsize{}
\end{flushright}

\problem{1}
Дана клетчатая фигура, покрашенная в шахматную раскраску. Известно, что в ней одинаковое количество клеток белого и чёрного цветов. Верно ли, что её можно разрезать на доминошки?

\problem{2}
(a) На шахматной доске $8 \times 8$ отмечены три клетки $A, B$ и $C$. В клетке $A$ стоит хромая ладья, которая за один ход умеет перемещаться в любую соседнюю по стороне клетку. Обязательно ли хромая ладья может обойти все клетки доски по одному разу, побывав сначала в клетке $B$, потом в клетке $C$ и закончив в клетке $A$?
(б) Из шахматной доски $8 \times 8$ вырезали 2 клетки разного цвета (в шахматной раскраске). Докажите, что оставшуюся часть можно разбить на доминошки.

\problem{3}
Из шахматной доски $8 \times 8$ вырезали 10 клеток.

(a) Какое наибольшее количество доминошек можно после этого гарантированно вырезать из этой доски?

(б) А если известно, что среди вырезанных клеток есть как чёрные, так и белые клетки?

\problem{4}
Можно ли доску $6 \times 6$ разбить на 18 доминошек так, чтобы любая прямая, параллельная сторонам квадрата, пересекала хотя бы одну доминошку?

\problem{5}
Клетчатый прямоугольник покрыт в два слоя доминошками (то есть каждая клетка принадлежит ровно двум доминошкам). Докажите, что доминошки можно разбить на два непересекающихся множества, каждое из которых покрывает весь прямоугольник.